\section{Commutation relations of the terminal matrix}
\label{rr-section}

We prove the full commutation relation of the actual terminal
matrix by adding one isotropic pair.  A full Capelli determinant
on the larger unitary group restricts to a product of three
$2\times2$ matrices.  The full determinants commute.
An elementary independence argument then recovers all sixteen
commutation relations of the middle matrix.
The smooth extension needed for this argument is given explicitly.

\subsection{The open cell with one isotropic pair}

Put $d=k+1$, let $J=\operatorname{diag}(1,\ldots,1,-1)$
on $W=\mathbb C^d$, and let
\[
 V=\mathbb C e\oplus W\oplus\mathbb C f,\qquad
 H=\begin{pmatrix}0&0&1\\0&J&0\\1&0&0\end{pmatrix}.
\]
For $v\in W$, $t\in\mathbb R$, and $c\in\mathbb C^\times$,
write
\[
 n(v,t)=
 \begin{pmatrix}
 1&-v^*J&it-\tfrac12v^*Jv\\
 0&I&v\\0&0&1
 \end{pmatrix},
 \qquad
 w(c)=
 \begin{pmatrix}0&0&\bar c^{-1}\\0&I&0\\c&0&0\end{pmatrix}.
\]
Both matrices preserve $H$.  The matrices $n(v,t)$ form
the unipotent group $N$.
For $z\in W$ set
\[
 \ell_z(n(v,t))=-v^*Jz,\qquad
              \chi_z(n)=\exp(\ell_z(n)-\overline{\ell_z(n)}).
\]
The coordinate $v$ is additive under multiplication in
$N$, so $\chi_z$ is a character.  This is the character
normalization used in the preceding radial calculations.

\begin{lemma}\label{rr-cell}
The map
\[
 N\times\mathbb C^\times\times U(W,J)\times N
   \longrightarrow \{g\in U(V,H):g_{fe}\ne0\},\qquad
 (n_L,c,m,n_R)\longmapsto n_Lw(c)m n_R
\]
is a smooth bijection with smooth inverse.
Here $m$ acts as the identity on $\mathbb Ce\oplus\mathbb Cf$.
\end{lemma}

\begin{proof}
Let $g$ be in the displayed open set and put $c=g_{fe}$.
Multiplication of the four displayed block matrices shows
that the only possible parameters satisfy
\[
 v_L=c^{-1}g_{We},\qquad
 v_R^*J=-c^{-1}g_{fW},\qquad
 a_L=c^{-1}g_{ee},\qquad a_R=c^{-1}g_{ff},
\]
where $a_\nu=it_\nu-\tfrac12v_\nu^*Jv_\nu$.
The zero Hermitian norm of column $e$ of $g$ gives
$2\operatorname{Re}a_L=-v_L^*Jv_L$.
The analogous row identity, obtained from $gHg^*=H$,
gives $2\operatorname{Re}a_R=-v_R^*Jv_R$.
Thus $t_L,t_R$ are uniquely the imaginary parts of
$a_L,a_R$.

Remove the two resulting unipotent factors.
The matrix $n_L^{-1}gn_R^{-1}$ has first column $cf$
and last row $ce^*$.
Its other columns are orthogonal to that first column,
so their $e$-coordinates vanish except in the last
column.  The analogous row orthogonality makes the
remaining off-diagonal blocks zero.  The pairing of
the first and last columns then fixes the $(e,f)$
entry to be $\bar c^{-1}$.
The remaining $W$-block $m$ satisfies $m^*Jm=J$.
It equals
$g_{WW}+v_Lc\,v_R^*J$ by multiplication of the factors.
This proves both existence and uniqueness.
All formulas use smooth operations and division only
by the nonzero number $c$, which proves smoothness.
\end{proof}

\begin{corollary}\label{rr-extension}
Every smooth function $\phi(c,m)$ on
$\mathbb C^\times\times U(W,J)$ extends uniquely to a
smooth function on this open cell with
\[
 F(n_Lw(c)m n_R)
                  =\chi_z(n_L)\chi_{z'}(n_R)\phi(c,m).
\]
It has the two prescribed unipotent character laws.
\end{corollary}

\begin{proof}
Use the coordinates of Lemma~\ref{rr-cell}.
Multiplication on either side by $N$ stays in the
cell, and its factor is multiplied by the corresponding
unipotent element.  The character law gives the assertion.
\end{proof}

\subsection{The full Capelli determinant in these coordinates}

Let $p=c\partial_c$, $\bar p=\bar c\partial_{\bar c}$
and $\gamma(c)=|c|^{d+1}$.  Thus
$p\gamma=\bar p\gamma=(d+1)\gamma/2$.
Define
\[
 T(Y)=\begin{pmatrix}Y+p&-c\\c^{-1}&0\end{pmatrix},
 \qquad
 L(Y)=\begin{pmatrix}Y-\bar p&\bar c^{-1}\\-\bar c&0\end{pmatrix}.
\]
Let $t(Y)$ be the actual terminal matrix from
Section~\ref{rq-section}.  Its entries act only on $m$
and commute with operators in $c$.

\begin{lemma}\label{rr-radial-capelli}
Let $\mathcal D(Y)$ be the full Capelli determinant on
$U(V,H)$ with parameter $Y+(d+1)/2$.
For every smooth function $F$ with the two unipotent
character laws, put
$\Phi(c,m)=\gamma(c)F(w(c)m)$.
Then
\begin{equation}\label{rr-sandwich}
 \gamma(c)(\mathcal D(Y)F)(w(c)m)
                         =(L(Y)t(Y)T(Y))_{11}\Phi(c,m).
\end{equation}
Consequently the operators
$\mathcal C(Y)=(L(Y)t(Y)T(Y))_{11}$ commute
for all complex parameters.
\end{lemma}

\begin{proof}
We give the bordered determinant calculation explicitly.
Order the full basis as $(e,1,\ldots,d,f)$ and put
$U=Y+(d+1)/2$.  For an omission in the middle block,
write $Q_{\widehat i,\widehat j}$ for its quantum minor.
Expansion first in column $e$, then in column $f$,
and then in the two exposed end rows gives the following
four terms:
\begin{align*}
 &T_{ee}(U)D(U-1)T_{ff}(U-d-1),\\
 &-\sum_{i,j}(-1)^{i+j}
      E_{fj}T_{ee}(U)Q_{\widehat i,\widehat j}(U-1)E_{if},\\
 &-\sum_{i,j}(-1)^{i+j}
      E_{ie}Q_{\widehat i,\widehat j}(U-2)
                        T_{ff}(U-d-1)E_{ej},\\
 &-\sum_{\substack{i\ne i'\\j\ne j'}}
   \epsilon(i,i')\epsilon(j,j')(-1)^{i+i'+j+j'}
   E_{ie}E_{fj}Q_{\widehat{i,i'},\widehat{j,j'}}(U-2)
                                      E_{ej'}E_{i'f}.
\end{align*}
The remaining terms are
\begin{align*}
 &(-1)^{d+1}E_{fe}Q_{(e,1,\ldots,d),(1,\ldots,d,f)}(U-1),\\
 &\sum_i(-1)^{i+d}E_{ie}
       Q_{(1,\ldots,\widehat i,\ldots,d,f),(1,\ldots,d)}(U-1)E_{ef}.
\end{align*}
These formulas are obtained using the end-row and
end-column expansions of Lemma~\ref{rq-cofactor-expansions}.
For the fourth line, the position of $i'$ after removing
$i$ is $i'-\mathbf1_{i'>i}$, and similarly for $j'$.
These position changes give the two factors $\epsilon$.
When an end-row factor is moved past an end-column
factor, their possible extra commutator is a multiple
of $E_{ef}$ or $E_{fe}$ and is exactly one of the last
two displayed terms.  Thus this expansion retains the
noncommutative correction terms.

At the point $w(c)m$, the left and right character
rules of \eqref{rc-mixed-right}--\eqref{rc-mixed-left}
are
\[
 \begin{gathered}
 E_{ej}F=x_jF,\qquad E_{if}F=\xi_iF,\\
 (E_{ie}XF)(w(c)m)=c^{-1}\eta_i(XF)(w(c)m),\\
 (E_{fj}XF)(w(c)m)=\bar c^{-1}y_j(XF)(w(c)m).
 \end{gathered}
\]
They hold for every right differential word $X$ in
the left rules.  The right central root $E_{ef}$
annihilates $F$, and the left central root $E_{fe}$
annihilates the evaluation of every $XF$.
The last two terms therefore vanish.

The diagonal derivatives restrict to $p$ and $-\bar p$.
Conjugation by $\gamma$ changes the two diagonal
factors to $Y+p$ and $Y-\bar p$, respectively.
The middle parameter $U-1$ is $Y+k/2$.
Grouping the four sums by their omitted pairs,
as in \eqref{rc-boundary-sign}, leaves precisely
\[
 (Y-\bar p)t_{11}(Y)(Y+p)
 +(Y-\bar p)t_{12}(Y)c^{-1}
 +\bar c^{-1}t_{21}(Y)(Y+p)
 +\bar c^{-1}t_{22}(Y)c^{-1}.
\]
This is the $(1,1)$ entry of $L(Y)t(Y)T(Y)$.

The full Capelli determinants commute, by their
centrality.  They preserve both unipotent character
laws: left translations commute with right
derivatives, and invariance under conjugation gives
commutation with right translations.  Corollary
\ref{rr-extension} allows an arbitrary smooth radial
function as their restriction.  Applying
\eqref{rr-sandwich} twice therefore proves
$[\mathcal C(U),\mathcal C(V)]=0$ on every such function.
\end{proof}

\subsection{An elementary matrix calculation}

For a $2\times2$ matrix $A$ over an algebra, let
$A_1=A\otimes I$ and $A_2=I\otimes A$, with their
entries ordered as written.  Let $P$ interchange the
two factors of $\mathbb C^2\otimes\mathbb C^2$ and put
$R(s)=sI-P$.
The identity
\[
 R(u-v)A_1(u)A_2(v)=A_2(v)A_1(u)R(u-v)
\]
is equivalent, entry by entry, to
\begin{equation}\label{rr-RTT}
 (u-v)[A_{ab}(u),A_{cd}(v)]
       =A_{cb}(u)A_{ad}(v)-A_{cb}(v)A_{ad}(u).
\end{equation}

\begin{lemma}\label{rr-Weyl}
If $cc^{-1}=c^{-1}c=1$ and $[p,c]=c$, then
\[
                  \begin{pmatrix}u+p&-c\\c^{-1}&0\end{pmatrix}
\]
satisfies \eqref{rr-RTT}.
Both $T$ and $L$ above therefore satisfy this identity.
\end{lemma}

\begin{proof}
One also has $[p,c^{-1}]=-c^{-1}$.
Commutators of two diagonal entries, or of $c$ with
$c^{-1}$, vanish.  The only nonzero pairs, up to
interchanging them, are
$[u+p,-c]=-c$ and $[u+p,c^{-1}]=-c^{-1}$.
Substitution into \eqref{rr-RTT} gives respectively
$-(u-v)c$ and $-(u-v)c^{-1}$ on both sides.
Pairs involving the zero entry give zero on both
sides.  This exhausts the four entries and proves
the identity.
For $L$, take the Weyl pair
$p'=-\bar p$, $c'=-\bar c^{-1}$,
$(c')^{-1}=-\bar c$.
\end{proof}

\begin{lemma}\label{rr-defect}
Let $Q(u)$ be any polynomial $2\times2$ operator
matrix acting only on $m$, and put $A(u)=L(u)Q(u)T(u)$.
Set
\[
 X(u,v)=R(u-v)Q_1(u)Q_2(v)
                         -Q_2(v)Q_1(u)R(u-v).
\]
Then
\begin{equation}\label{rr-projected-defect}
 (u-v-1)[A_{11}(u),A_{11}(v)]
  =\bigl(L_2(v)L_1(u)X(u,v)T_1(u)T_2(v)\bigr)_{11,11}.
\end{equation}
\end{lemma}

\begin{proof}
The two Weyl pairs commute with one another and with
all entries of $Q$.  Move $R$ through the two $L$
factors by Lemma~\ref{rr-Weyl}, and then through the
two $T$ factors by the same identity.  Ordinary
associative matrix multiplication gives
\[
 R A_1(u)A_2(v)-A_2(v)A_1(u)R
                       =L_2(v)L_1(u)X(u,v)T_1(u)T_2(v).
\]
The row and column indexed by $(1,1)$ are eigenvectors
of $P$ with eigenvalue one.  Taking that entry gives
\eqref{rr-projected-defect}.
\end{proof}

\subsection{Recovering every coefficient}

\begin{lemma}\label{rr-independence}
Suppose the left side of
\eqref{rr-projected-defect} vanishes on all smooth
functions of $(c,m)$.  If $u-v\ne1$, then
$X(u,v)$ is the zero operator matrix on smooth
functions of $m$.
\end{lemma}

\begin{proof}
Use indices $0,1$ in this calculation only.
Fix a smooth function $f(m)$, a point $m$, a complex
number $\beta$ and an integer $n$, and apply
\eqref{rr-projected-defect} to
\[
                     |c|^{2\beta}c^n f(m).
\]
Here $|c|^{2\beta}=\exp(2\beta\log|c|)$ is a smooth
single-valued function on $\mathbb C^\times$.
It is an eigenfunction of $p$ with eigenvalue
$\beta+n$ and of $\bar p$ with eigenvalue $\beta$.
Put $z=\beta+n$.

For $A=(a_1,a_2)$ and $B=(b_1,b_2)$ the first-row
and first-column coefficients in the right side of
\eqref{rr-projected-defect}, after removing the
common nonzero factor $|c|^{2\beta}c^n$, are
\[
      L_A(\beta)R_B(z)c^{-(b_1+b_2)}\bar c^{-(a_1+a_2)},
\]
where
\[
\begin{array}{c|cccc}
A&00&01&10&11\\ \hline
L_A(\beta)&(u-\beta)(v-\beta)&u-\beta&v-\beta+1&1
\end{array}
\]
and
\[
\begin{array}{c|cccc}
B&00&01&10&11\\ \hline
R_B(z)&(u+z)(v+z)&u+z-1&v+z&1.
\end{array}
\]
The two shifts by one follow from
$p(c^{-1}h)=c^{-1}(p-1)h$ and
$\bar p(\bar c^{-1}h)=\bar c^{-1}(\bar p-1)h$.
No entry of $X$ differentiates these scalar factors.

The nine monomials
$c^{-b}\bar c^{-a}$, $0\leq a,b\leq2$,
are linearly independent on every nonempty open subset
of $\mathbb C^\times$.  Indeed multiplying by
$c^2\bar c^2$ gives a polynomial in $c,\bar c$;
the invertible change $c=x+iy$, $\bar c=x-iy$
reduces the assertion to polynomial uniqueness on
an open subset of $\mathbb R^2$.
Thus for each pair $(a,b)$,
\[
 \sum_{\substack{a_1+a_2=a\\b_1+b_2=b}}
 L_A(\beta)R_B(\beta+n)(X_{A,B}(u,v)f)(m)=0.
\]
This holds for every complex $\beta$ and every
integer $n$.  Polynomial uniqueness first in
$\beta$ and then in $n$ makes it a polynomial
identity in two independent variables $\beta,z$.
The argument remains valid for vector-valued
coefficients by finite Vandermonde inversion.

When $a=0$ or $2$, there is just one nonzero row
polynomial.  When $a=1$, the two row polynomials
are $u-\beta$ and $v-\beta+1$.
They are independent because their difference is
the nonzero constant $u-v-1$.
The same assertion holds for the two column
polynomials $u+z-1$ and $v+z$ when $b=1$.
Products of these independent polynomials in the
two separate variables are independent.
Every coefficient $(X_{A,B}(u,v)f)(m)$ therefore
vanishes.  Since $f,m$ were arbitrary, $X(u,v)=0$.
\end{proof}

\begin{theorem}\label{rr-terminal-RTT}
The actual terminal differential matrix satisfies,
on every smooth function on $U(W,J)$,
\[
 (u-v)[t_{ab}(u),t_{cd}(v)]
      =t_{cb}(u)t_{ad}(v)-t_{cb}(v)t_{ad}(u)
                \qquad(a,b,c,d\in\{1,2\}).
\]
The statement holds for every pair of complex
parameters, including $u-v=1$.
\end{theorem}

\begin{proof}
Lemma~\ref{rr-radial-capelli} gives the commuting
$(1,1)$ entries required by Lemma~\ref{rr-independence}.
Hence the defect matrix vanishes for $u-v\ne1$.
All its entries are polynomials in $u,v$ with
differential-operator coefficients, since the
terminal entries are finite sums of quantum minors.
Applying an entry to any fixed smooth function and
evaluating at a point gives a polynomial which
vanishes for $u-v\ne1$; it therefore vanishes
identically.  This includes the exceptional parameters
and proves the theorem.
\end{proof}

\begin{corollary}\label{rr-terminal-adjugate}
Put $J_2=\begin{pmatrix}0&1\\-1&0\end{pmatrix}$ and
\[
 q(Y)=t_{11}(Y+1)t_{22}(Y)-t_{21}(Y+1)t_{12}(Y).
\]
As operator matrices on smooth terminal functions,
\[
 t(Y+1)^{\mathsf T}J_2t(Y)=q(Y)J_2,\qquad
 t(Y)J_2t(Y+1)^{\mathsf T}=q(Y)J_2.
\]
On the compact types of Corollary~\ref{rq-type-scalar},
$q(Y)$ is the scalar polynomial computed there.
At a value where that scalar is nonzero,
\[
             t(Y)^{-1}=-q(Y)^{-1}
                              J_2t(Y+1)^{\mathsf T}J_2
\]
is both a left and a right inverse.
\end{corollary}

\begin{proof}
Write $S=t(Y+1)$ and $T=t(Y)$.
Use Theorem~\ref{rr-terminal-RTT} with parameter
difference one.  It gives
\[
 S_{11}T_{21}=S_{21}T_{11},\qquad
 S_{12}T_{22}=S_{22}T_{12},
\]
and, on adding its $(12,21)$ and $(11,22)$ entries,
\[
 S_{12}T_{21}-S_{22}T_{11}
                       =-S_{11}T_{22}+S_{21}T_{12}.
\]
These are the two zero diagonal entries and the
negative lower off-diagonal entry of the first
matrix identity; its upper off-diagonal entry
is the definition of $q(Y)$.

The same commutation relation, in its $(11,12)$,
$(21,12)$, $(11,22)$, and $(21,22)$ entries, gives
\begin{align*}
 T_{11}S_{12}-T_{12}S_{11}&=0,\\
 T_{11}S_{22}-T_{12}S_{21}&=q(Y),\\
 T_{21}S_{12}-T_{22}S_{11}&=-q(Y),\\
 T_{21}S_{22}-T_{22}S_{21}&=0.
\end{align*}
This proves the second matrix identity.
Apply Corollary~\ref{rq-type-scalar} for the scalar
value.  Finally $J_2^2=-I$, so the two identities
give the displayed two-sided inverse whenever
that scalar is nonzero.
\end{proof}

\paragraph{Scope.}
The theorem proves the full commutation relation of
the terminal matrix itself, without assuming a
finite compact type or a special Jacquet origin
for the smooth function.  The recursive frame
identities and the estimates justifying integration
with the Jacquet pairing remain separate.


